% =====================================================================
\section{Descrição Formal do Mundo dos Blocos de Tamanho Variado}
\label{sec:formal}
% =====================================================================
\responsavel{(nome)}

\subsection{Domínio}
\afazer{Blocos e comprimentos, pontos e slots, níveis, instantes, $\spn(b,p)$.}

\subsection{Predicados e axiomas}
\afazer{Tabela dos predicados ($\at$, $\lev$, $\cov$, $\livre$, $\esp$, $\clr$, $\on$) e
os axiomas em LPO (unicidade, exclusão, estabilidade etc.).}

\subsection{A ação \texorpdfstring{$\mv(b,y,p)$}{move(b,y,p)}}
\afazer{Pré-condições (P1, P2, \dots) e efeitos em estilo STRIPS (ADD / DEL).}

\subsection{Elementos novos e as ações associadas}
\afazer{Item 2 do enunciado: para cada elemento novo, o que a ação adiciona e remove.}

\subsection{Ajustes em relação ao manual}
\afazer{O que foi mudado em relação ao manual e por quê (predicado $\esp$ no lugar de
$\clr(y)$; regra da sombra).}

\subsection{Execução manual}
\label{sec:manual}
\subsubsection{Situação 1}
\responsavel{(nome)}
\afazer{Figura com os estados do plano e tabela: ação, pré-condições, ADD, DEL.
Argumento de por que o plano é mínimo.}

\subsubsection{Situação 2}
\responsavel{(nome)}
\afazer{Idem para $S_0 \to S_5$.}

\subsubsection{Situação 3}
\responsavel{(nome)}
\afazer{Idem para $S_0 \to S_7$.}

\subsection{Ordem parcial}
\label{sec:pop}
\afazer{Definição: ações, vínculos causais $A\xrightarrow{\;p\;}B$, restrições de
ordem, ameaças.}

\subsubsection{Situação 1}
\afazer{Tabela de vínculos causais, grafo, análise de ameaças, linearizações.}

\subsubsection{Situação 2}
\afazer{Idem.}

\subsubsection{Situação 3}
\afazer{Idem.}

\subsubsection{Ordem parcial no SAT}
\afazer{Como as restrições $\varphi_1\prec_P\varphi_2$ viram cláusulas e o resultado
de cada cenário.}
